\chapter{Doxometric Measurement}
\label{ch:doxometric-measurement}

\epigraph{Not everything that counts can be counted, and not everything that can be counted counts.}{William Bruce Cameron}

\noindent
Doxonomics requires measurement.
Without it, the field remains assertion---a collection of plausible claims about belief production that cannot be tested, compared, or refuted.
This chapter develops the theory of \emph{doxometrics}\index{doxometrics}: the quantitative measurement of public belief, its intensity, distribution, and institutional embeddedness.

The central difficulty is that beliefs are \emph{latent variables}.
We do not observe them directly.
We observe indicators---survey responses, behaviors, choices, texts, institutional practices---and must infer the underlying belief structure.
This inference is never perfect, and the measurement challenges are at least as severe as those in any social science.
Doxometrics must therefore be psychometrically serious: attentive to reliability, validity, bias, and the gap between what we measure and what we claim to measure.


\section{What Counts as a Measurable Belief?}
\label{sec:ch13-operationalization}

\begin{definition}[Operationalized Belief]
\label{def:operationalized}
\index{operationalized belief}
A belief $p$ is \emph{doxometrically operationalized} when there exists a measurement function $m: \text{Pop} \to [0,1]$ such that $m(j)$ approximates $\pi_j(p)$ with known reliability and validity.
The population-level measure is $\hat{\mu}(p) = \frac{1}{N}\sum_{j=1}^N m(j)$.
\end{definition}

Not all beliefs are equally measurable.
The layered ontology (Principle~\ref{prin:layered-ontology}) implies that different layers present different measurement challenges:

\begin{center}
\begin{tabular}{lp{4.5cm}p{5cm}}
\toprule
\textbf{Layer} & \textbf{Measurement method} & \textbf{Difficulty} \\
\midrule
1. Propositional belief & Surveys, polls & Moderate (social desirability bias, framing effects) \\
2. Norm perception & Norm-elicitation surveys, vignettes & Moderate (perception $\neq$ reality) \\
3. Practical schema & Behavioral observation, reaction time & High (must be inferred from action, not report) \\
4. Identity commitment & Identity scales, in-group/out-group tasks & Moderate-high (self-report may be performative) \\
5. Affective orientation & Implicit association tests, physiological measures & High (requires specialized equipment; validity debated) \\
\bottomrule
\end{tabular}
\end{center}

Most doxometric work focuses on layer~1 because it is the most tractable.
But the most consequential doxonomic products---deep anthropological regimes, embodied dispositions, affective prejudices---operate at layers 3--5, where measurement is hardest.
This gap between what is most measurable and what is most important is the field's central methodological challenge.


\section{Indicators and Proxies}
\label{sec:ch13-indicators}

Direct measurement of $\pi_j(p)$ is rarely possible.
We distinguish four types of indicators, each with characteristic strengths and biases:

\begin{description}[nosep]
  \item[Stated belief:] direct survey response (``Do you agree that human nature is selfish?'').
  \textit{Strength}: easy to collect at scale.
  \textit{Bias}: social desirability (respondents give ``acceptable'' answers), acquiescence (tendency to agree with any statement), framing effects (wording changes responses dramatically).

  \item[Behavioral marker:] observed action consistent with $p$ (e.g., defection rates in public goods games as a proxy for selfishness belief; willingness to cross a picket line as a proxy for individualist values).
  \textit{Strength}: resistant to social desirability bias (actions are costly, so they reveal more than words).
  \textit{Bias}: behavior is multiply determined (defection in a game may reflect strategic calculation, not belief about human nature).

  \item[Revealed preference:] choices that imply $p$ (e.g., voting for policies consistent with meritocratic belief; purchasing choices consistent with consumerist identity).
  \textit{Strength}: based on high-stakes real-world decisions.
  \textit{Bias}: preferences are not beliefs (I may vote for redistribution because it benefits me, not because I reject meritocracy).

  \item[Textual indicator:] frequency and framing of $p$ in media, social media, institutional documents, or personal communication.
  \textit{Strength}: available at massive scale through computational methods.
  \textit{Bias}: discourse frequency is not belief prevalence (a proposition may be frequently discussed precisely because it is contested, not because it is believed).
\end{description}

\begin{proposition}[Divergence of Indicators]
\label{prop:divergence}
Stated belief and behavioral markers may diverge systematically.
A population may report cooperative values while behaving competitively (or vice versa).
The gap $|m_{\text{stated}}(p) - m_{\text{behavioral}}(p)|$ is itself a doxonomic observable, measuring the \emph{depth of internalization}.
\end{proposition}

\begin{remark}[The Divergence as Data]
\label{rem:divergence-as-data}
When stated belief exceeds behavioral compliance, the belief has not been internalized to layer~3 (practical schema): people profess it but do not enact it.
When behavioral compliance exceeds stated belief, the belief may operate doxically (layer~5): people enact it without explicitly endorsing it.
The gap between layers is not noise to be eliminated; it is signal about the depth and mode of doxonomic penetration.
\end{remark}

\begin{example}[Measuring Belief in Meritocracy]
\label{ex:measuring-meritocracy}
Four indicators for the belief ``success reflects effort and talent'':
\begin{enumerate}[nosep]
  \item \textbf{Stated}: Likert-scale agreement with ``people who work hard almost always get ahead'' (World Values Survey item).
  \item \textbf{Behavioral}: willingness to assign unequal rewards in an experimental dictator game after observing performance differences.
  \item \textbf{Revealed preference}: voting for politicians who oppose redistributive taxation.
  \item \textbf{Textual}: frequency of meritocratic framing (``earned,'' ``deserved,'' ``self-made'') in media coverage of wealth.
\end{enumerate}
Each captures a different facet.
A complete doxometric assessment uses all four and reports both the individual estimates and the divergences.
\end{example}


\section{Measurement Bias and Correction}
\label{sec:ch13-bias}

Several systematic biases plague doxometric measurement.

\begin{definition}[Social Desirability Bias]
\label{def:social-desirability}
\index{social desirability bias}
\emph{Social desirability bias} (SDB) is the tendency of respondents to answer survey questions in a way that will be viewed favorably by others, rather than truthfully.
In doxometric contexts, SDB causes under-reporting of socially disapproved beliefs (racism, selfish motivation) and over-reporting of socially approved beliefs (cooperation, fairness).
\end{definition}

Correction strategies:
\begin{itemize}[nosep]
  \item \textbf{List experiments}\index{list experiment}: respondents report the number of statements they agree with from a list, rather than identifying which ones.
  The treatment group's list includes the sensitive item; the control group's does not.
  The difference in average counts estimates the prevalence of the sensitive belief without any individual revealing their answer.
  \item \textbf{Endorsement experiments}: respondents evaluate a policy endorsed by a figure associated with the belief.
  The endorsement effect measures implicit alignment without direct self-report.
  \item \textbf{Implicit measures}: reaction-time tasks (Implicit Association Test) measure automatic associations, bypassing deliberate self-presentation.
  \textit{Caveat}: the reliability and validity of IATs for measuring beliefs (as opposed to associations) remains debated.
  \item \textbf{Behavioral validation}: compare survey responses with behavior in incentivized tasks.
  Respondents who claim to value cooperation but defect in public goods games are flagged as exhibiting SDB.
\end{itemize}

\begin{definition}[Framing Bias]
\label{def:framing-bias}
\index{framing bias}
\emph{Framing bias} occurs when the wording of a survey question influences the response, independent of the underlying belief.
``Do you believe in free enterprise?'' and ``Do you think corporations should be unregulated?'' may target the same proposition but elicit different responses.
Doxometric instruments should include multiple framings of the same target belief and report the framing sensitivity as a measure of measurement precision.
\end{definition}

\begin{definition}[Measurement Invariance]
\label{def:invariance}
\index{measurement invariance}
A doxometric instrument satisfies \emph{measurement invariance} across groups $G_1, \ldots, G_L$ if the relationship between latent belief and observed indicators is the same in each group.
Formally: the factor loadings $\lambda_k$ and intercepts are equal across groups (metric and scalar invariance).
Without invariance, cross-group comparisons of belief levels are not meaningful---a higher score in group $G_1$ might reflect a stronger belief or a different relationship between the belief and the indicator.
\end{definition}

\begin{remark}[Why Invariance Matters for Doxonomics]
\label{rem:invariance-matters}
Doxonomic analysis frequently compares belief levels across countries, classes, racial groups, or media environments.
If the measurement instrument is not invariant across these groups---if ``hard work leads to success'' means something different to a white professional in Connecticut than to a Black factory worker in Detroit---then the comparison is spurious.
Every cross-group doxometric comparison should be accompanied by an invariance test.
\end{remark}


\section{Latent Variable Models}
\label{sec:ch13-latent}

\begin{model}[Factor Model for Belief]
\label{mod:factor}
\index{latent variable model}
Let $\bm{y}_j = (y_{j1}, \ldots, y_{jK})$ be $K$ observable indicators for agent $j$.
The latent belief $\pi_j$ generates observables via
\[
  y_{jk} = \lambda_k \pi_j + \varepsilon_{jk}
\]
where $\lambda_k$ is the factor loading for indicator $k$ (how strongly the indicator reflects the latent belief) and $\varepsilon_{jk} \sim \mathcal{N}(0, \sigma_k^2)$ is measurement error.
The latent belief is recovered as the weighted average:
\[
  \hat{\pi}_j = \frac{\sum_k \lambda_k y_{jk} / \sigma_k^2}{\sum_k \lambda_k^2 / \sigma_k^2}
\]
This is the minimum-variance unbiased estimator under the model.
\end{model}

\begin{proposition}[Reliability]
\label{prop:reliability}
\index{reliability}
The \emph{reliability} of the latent belief estimate is
\[
  \text{Rel}(\hat{\pi}) = 1 - \frac{\mathrm{Var}(\hat{\pi} - \pi)}{\mathrm{Var}(\hat{\pi})} = \frac{\sum_k \lambda_k^2 / \sigma_k^2}{\sum_k \lambda_k^2 / \sigma_k^2 + 1}
\]
Reliability increases with the number of indicators $K$ (more items → more signal) and with the factor loadings $\lambda_k$ (stronger items → more signal), and decreases with measurement error variance $\sigma_k^2$.
The standard threshold for acceptable reliability in psychometrics is $\text{Rel} > 0.70$.
\end{proposition}

\begin{remark}[Beyond Single-Factor Models]
\label{rem:multi-factor}
The single-factor model assumes that all indicators load on a single latent belief.
Many doxonomic beliefs are multi-dimensional.
``Meritocratic belief'' may decompose into: belief in effort → success, belief in equal opportunity, belief in desert-based distribution, and attribution of poverty to individual failure.
These may be correlated but not identical.
Exploratory factor analysis (EFA) and confirmatory factor analysis (CFA) should be used to assess dimensionality before constructing composite indices.
Treating a multi-dimensional construct as unidimensional produces unreliable and uninterpretable results.
\end{remark}

\subsection{Item Response Theory}

For more refined measurement, item response theory (IRT) models the probability of endorsing each item as a function of the latent belief and item characteristics:

\begin{model}[Two-Parameter IRT]
\label{mod:irt}
\index{item response theory}
The probability that agent $j$ endorses item $k$ is
\[
  P(y_{jk} = 1 \mid \pi_j) = \frac{1}{1 + \exp(-a_k(\pi_j - b_k))}
\]
where $a_k$ is the \emph{discrimination parameter} (how sharply the item distinguishes between high and low $\pi_j$) and $b_k$ is the \emph{difficulty parameter} (the belief level at which endorsement probability is 50\%).
\end{model}

IRT improves on the factor model by allowing each item to have a different measurement precision at different points on the belief continuum.
A well-designed doxometric instrument should include items with different difficulty parameters to measure accurately across the full range of belief levels---from strong disbelief through ambivalence to strong endorsement.


\section{Composite Indices of Belief Dominance}
\label{sec:ch13-composite}

\begin{definition}[Doxonomic Dominance Index]
\label{def:dominance-index}
\index{dominance index}
The \emph{doxonomic dominance index} (DDI) for belief $p$ is
\[
  \mathrm{DDI}(p) = \alpha_1 \hat{\mu}(p) + \alpha_2 \cdot (1 - \mathrm{Var}(\hat{\pi})) + \alpha_3 \cdot \mathrm{InstEmbed}(p) + \alpha_4 \cdot \mathrm{MediaFreq}(p)
\]
where:
\begin{itemize}[nosep]
  \item $\hat{\mu}(p)$ is the population-mean belief (how widely held),
  \item $1 - \mathrm{Var}(\hat{\pi})$ is the concentration of belief (how uniformly held),
  \item $\mathrm{InstEmbed}(p)$ measures institutional embeddedness (how many institutions treat $p$ as operational truth),
  \item $\mathrm{MediaFreq}(p)$ is the normalized frequency of $p$ in public discourse.
\end{itemize}
The weights $\alpha_i$ satisfy $\sum \alpha_i = 1$.
\end{definition}

\begin{remark}[DDI Limitations]
\label{rem:ddi-limitations}
The DDI is a \emph{dashboard sketch}, not a validated index.
The weights $\alpha_i$ are currently arbitrary; different weights produce different rankings.
The components mix unlike quantities (survey proportions, variance measures, frequency counts) without a clear scaling doctrine.
Before the DDI can serve as a serious research tool, it requires:
\begin{itemize}[nosep]
  \item empirical validation (does DDI predict policy outcomes, behavioral patterns, or other independently measurable consequences?),
  \item sensitivity analysis (how much do results change with different weight specifications?),
  \item normalization (each component should be on a comparable scale).
\end{itemize}
We include it here as a conceptual prototype, not as a finished instrument.
\end{remark}


\section{Measuring Institutional Embeddedness}
\label{sec:ch13-embeddedness}

The DDI requires a measure of institutional embeddedness---how deeply a belief is woven into the operational assumptions of institutions.

\begin{definition}[Institutional Embeddedness]
\label{def:inst-embed}
\index{institutional embeddedness}
The \emph{institutional embeddedness} of belief $p$ is
\[
  \mathrm{InstEmbed}(p) = \frac{1}{|\mathcal{U}|} \sum_{U_i \in \mathcal{U}} \mathbf{1}[p \in \sigma_i]
\]
where $\mathcal{U}$ is the set of institutional contexts and $\mathbf{1}[p \in \sigma_i] = 1$ if $p$ functions as an operational assumption in context $U_i$.
\end{definition}

Operationalizing this requires coding institutional practices, curricula, regulations, and organizational structures for the presence of $p$ as an operational assumption---a qualitative judgment that can be systematized through:
\begin{itemize}[nosep]
  \item textbook analysis (does the economics curriculum assume rational self-interest?),
  \item organizational document analysis (do HR policies assume shirking? do compensation structures assume competitive motivation?),
  \item legal analysis (does regulatory framework assume market efficiency?),
  \item expert survey (do practitioners in a field report that $p$ is a taken-for-granted assumption in their work?).
\end{itemize}

\begin{example}[Embeddedness of ``Rational Self-Interest'']
\label{ex:embeddedness}
Coding for the belief ``agents are rational self-interest maximizers'' across U.S.\ institutional contexts:
\begin{center}
\begin{tabular}{lcc}
\toprule
\textbf{Institution} & \textbf{Embedded?} & \textbf{Evidence} \\
\midrule
Economics curricula & Yes & Textbook analysis \\
Corporate HR / compensation & Yes & Incentive-based pay design \\
Antitrust regulation & Yes & Consumer welfare standard \\
Criminal justice & Partial & Deterrence theory; also retributive \\
Public education & No & Assumes cooperation, development \\
Religious institutions & No & Assumes altruism, duty \\
Military & Partial & Hierarchy + unit solidarity \\
Family law & No & Assumes care, obligation \\
\bottomrule
\end{tabular}
\end{center}
$\mathrm{InstEmbed} \approx 3.5/8 = 0.44$.
The belief is institutionally embedded in roughly half of major contexts---far from total, but substantial enough to function as a default assumption in public life.
\end{example}


\section{Measuring Doxonomic Change Over Time}
\label{sec:ch13-temporal}

Doxonomic analysis requires not just snapshot measurement but the tracking of belief trajectories over time.

\begin{definition}[Belief Trajectory]
\label{def:trajectory}
\index{belief trajectory}
The \emph{belief trajectory} for proposition $p$ is the time series $\{\hat{\mu}(p, t)\}_{t=0}^{T}$.
A doxonomic event (a major campaign, a crisis, a curriculum reform) produces a detectable shift in the trajectory.
\end{definition}

Sources for longitudinal belief data:
\begin{itemize}[nosep]
  \item \textbf{Repeated cross-sections}: World Values Survey (waves every 5 years since 1981), General Social Survey (annual/biennial since 1972), Eurobarometer (biannual since 1973).
  \item \textbf{Panel surveys}: tracking the same individuals over time (more informative but more expensive; subject to attrition).
  \item \textbf{Text time series}: media corpus analysis tracking narrative frequency and framing over decades (Google Ngrams, GDELT, Media Cloud).
  \item \textbf{Behavioral time series}: cooperation rates in experimental games, charitable giving rates, labor market indicators.
\end{itemize}

\begin{remark}[Identifying Doxonomic Events in Time Series]
\label{rem:event-identification}
A sudden shift in $\hat{\mu}(p, t)$ coinciding with a known doxonomic investment (a new think tank, a major ad campaign, a curriculum reform, a media ownership change) is \emph{suggestive} but not \emph{causal} evidence of doxonomic effect.
Establishing causality requires the methods of Chapter~\ref{ch:causal-inference}: natural experiments, difference-in-differences, or synthetic controls that rule out confounding events.
Time-series correlation between funding and belief is a starting point, not a conclusion.
\end{remark}


\section{Notes and References}
\label{sec:ch13-notes}

Measurement theory for social science draws on \textcite{king1994} and \textcite{gelman2020}.
Latent variable models and factor analysis are standard in psychometrics; see Bollen (1989) for the structural equation modeling framework.
Item response theory is developed in Embretson and Reise (2000).
Behavioral measures of beliefs about cooperation are developed in \textcite{fischbacher2001} and \textcite{henrich2001}.
Survey methodology for the digital age connects to \textcite{salganik2018}.
On the divergence between stated and revealed preferences, see \textcite{kahneman2011}.

List experiments and other techniques for sensitive questions are reviewed in Blair, Imai, and Zhou (2015).
The Implicit Association Test is from Greenwald, McGhee, and Schwartz (1998); for a critical assessment, see Oswald et al.\ (2013).
Measurement invariance testing is covered in Meredith (1993) and Vandenberg and Lance (2000).
For topic models as tools for measuring latent constructs in text, see \textcite{blei2003} and \textcite{grimmer2013}.
On the distinction between discourse frequency and belief prevalence, see \textcite{zaller1992}: media coverage shapes the \emph{accessibility} of considerations, not directly the \emph{beliefs} of the audience.


\section{Exercises}

\begin{exercise}
Design a five-item survey instrument to measure belief in selfish human nature.
Include at least two reverse-coded items.
Specify the response format, argue for content validity, and describe how you would assess reliability and dimensionality.
\end{exercise}

\begin{exercise}
Using the factor model (Model~\ref{mod:factor}), show that the reliability of the latent belief estimate $\hat{\pi}_j$ increases with the number of indicators $K$ and decreases with measurement error variance $\sigma_k^2$.
Compute the minimum $K$ needed for $\text{Rel} > 0.80$ when $\lambda_k = 0.7$ and $\sigma_k = 0.5$ for all $k$.
\end{exercise}

\begin{exercise}
Compute the DDI for the belief ``markets are efficient'' using plausible estimates for each component.
Then repeat with different weight specifications ($\alpha_1 = 0.5$ vs.\ $\alpha_1 = 0.1$).
How sensitive is the DDI to the choice of weights?
What does this imply about its usefulness?
\end{exercise}

\begin{exercise}
\label{ex:ch13-list}
Design a list experiment to measure the prevalence of the belief ``most poor people are poor because they are lazy'' in a population where this belief is socially disapproved.
Specify the control and treatment lists, the sample size calculation, and the estimator.
\end{exercise}

\begin{exercise}
\label{ex:ch13-invariance}
A doxometric instrument measuring meritocratic belief is administered in the U.S.\ and Sweden.
Describe how you would test for measurement invariance across the two countries.
What would you conclude if configural invariance holds but scalar invariance fails?
Can you compare mean belief levels across the two countries?
\end{exercise}

\begin{exercise}
\label{ex:ch13-layers}
Choose a specific belief (e.g., ``inequality is natural'').
For each of the five layers in the ontology (Principle~\ref{prin:layered-ontology}), propose a concrete measurement method.
Which layers are most important for doxonomic analysis?
Which are most measurable?
What is the gap between importance and measurability, and what does it imply for the field?
\end{exercise}

\begin{exercise}
Using publicly available World Values Survey data, plot the trajectory of belief in ``hard work leads to success'' in three countries over the past 30 years.
Do you observe any doxonomic events (sudden shifts)?
If so, what was happening in each country at the time of the shift?
\end{exercise}

\begin{exercise}
The chapter notes that discourse frequency is not the same as belief prevalence.
Give an example where $\mathrm{MediaFreq}(p)$ is high but $\hat{\mu}(p)$ is low, and vice versa.
What does each configuration imply doxonomically?
\end{exercise}
